\chapter{Migrating to Bijou 7.0.0}\label{record-270}

\noindent{\footnotesize\color{BookMuted}\path{docs/releases/7.0.0/migration-guide.md}\par}\smallskip

Bijou 7.0.0 carries the V6 layout floor and V7 Product Truth lane into one major release. Most application code that stayed on public APIs should keep working, but maintainers should check the areas below before upgrading.

\section{Keep The 5.0 Hosted App Entrypoints}\label{record-270-keep-the-5.0-hosted-app-entrypoints}

The preferred hosted app path from 5.0.0 remains the right path:

\begin{CodeBlock}
const app = createFramedApp(options);
await app.run({ ctx });
\end{CodeBlock}

or:

\begin{CodeBlock}
await runFramedApp(options, { ctx });
\end{CodeBlock}

If your app still calls lower-level runners directly, this is a good release to move back to the frame-owned hosted entrypoints so shell timing, theme, and runtime behavior stay in one place.

\section{Check Table Width Expectations}\label{record-270-check-table-width-expectations}

\texttt{t\allowbreak{}a\allowbreak{}b\allowbreak{}l\allowbreak{}e\allowbreak{}(\allowbreak{})} and \texttt{t\allowbreak{}a\allowbreak{}b\allowbreak{}l\allowbreak{}e\allowbreak{}S\allowbreak{}u\allowbreak{}r\allowbreak{}f\allowbreak{}a\allowbreak{}c\allowbreak{}e\allowbreak{}(\allowbreak{})} now use bounded-width behavior by default for human-mode output. If your app expected tables to grow to intrinsic content width regardless of terminal size, opt into intrinsic layout explicitly.

For most terminal UIs, the new default is what you want: columns fit the target width, wrapped cells preserve word boundaries when possible, and lower-mode pipe formats remain deterministic.

\section{Prefer Standard Blocks For Product Semantics}\label{record-270-prefer-standard-blocks-for-product-semantics}

Bijou 7.0.0 exports a much broader standard Block catalog. If your app has custom semantic surfaces for status, feedback, navigation, disclosure, dense comparison, hierarchy, text entry, or documentation-like content, evaluate the new standard Blocks before inventing another local schema.

Blocks should own product semantics, data contracts, command intents, and lowering facts. Leaf components should remain the rendering vocabulary.

\section{Route Localization Through The Runtime Port}\label{record-270-route-localization-through-the-runtime-port}

\texttt{@\allowbreak{}f\allowbreak{}l\allowbreak{}y\allowbreak{}i\allowbreak{}n\allowbreak{}g\allowbreak{}r\allowbreak{}o\allowbreak{}b\allowbreak{}o\allowbreak{}t\allowbreak{}s\allowbreak{}/\allowbreak{}b\allowbreak{}i\allowbreak{}j\allowbreak{}o\allowbreak{}u\allowbreak{}-\allowbreak{}i\allowbreak{}1\allowbreak{}8\allowbreak{}n} now provides structured localization results instead of only returning strings. New localization code should preserve the returned metadata: key, locale, direction, entry kind, fallback or missing status, issues, and facts.

If your app persists locale preference, keep that persistence in a host adapter or other boundary-owned port. DOGFOOD uses this shape so runtime locale changes are visible without making app logic read process globals directly.

\section{Notification Overlays Need No Workaround}\label{record-270-notification-overlays-need-no-workaround}

Live notification rows now wrap at word boundaries when possible. Remove any app-local workaround that pre-wrapped notification titles, bodies, or action labels only to avoid split words in Bijou overlays.

Long unbroken text still hard-wraps inside the assigned notification width, so callers do not need to pre-truncate machine identifiers or hashes for layout safety.

\section{Release And DOGFOOD Proof Are Stricter}\label{record-270-release-and-dogfood-proof-are-stricter}

DOGFOOD now reads version-specific release docs from:

\begin{CodeBlock}
docs/releases/<package-version>/whats-new.md
docs/releases/<package-version>/migration-guide.md
\end{CodeBlock}

When preparing a release branch, bumping package versions without adding those docs will fail the docs and DOGFOOD test lanes. This is intentional: the package version and current release documentation must move together.

\section{Compatibility Stance}\label{record-270-compatibility-stance}

Bijou 7.0.0 is a major release because public contracts are broader and more explicit, not because every app must rewrite. The main migration work is to remove local workarounds that are now framework-owned and to adopt the standard contracts where they replace app-local semantic surfaces.

